\chapter{Geometry and Three-Manifolds}

\begin{definition}[Knot and link]
\label{thurston-ch01-knot-link}
\group{thurston-ch01-topology}
\source{thurston:page-0009}
\notready
A knot in $S^3$ is a closed simple curve, and a link is a finite union of pairwise disjoint closed simple curves.  Removing the interior of a tubular neighborhood produces a compact manifold with torus boundary.
\end{definition}

\begin{proposition}[Knot-complement fundamental group]
\label{thurston-ch01-knot-complement-group}
\group{thurston-ch01-topology}
\source{thurston:page-0010}
\notready
The fundamental group of a knot complement is isomorphic to $\mathbb Z$ exactly when the knot is trivial.
\end{proposition}
\begin{proof}
\sketch
The trivial knot has complement with fundamental group $\mathbb Z$; the converse is the Papakyriakopoulos theorem.
\end{proof}

\begin{definition}[Dehn surgery]
\label{thurston-ch01-dehn-surgery}
\group{thurston-ch01-topology}
\source{thurston:page-0010}
\notready
Dehn surgery removes a tubular neighborhood of every component of a knot or link and glues one solid torus to each resulting boundary torus.  On each boundary component, the gluing is determined by the primitive nonzero element of $\pi_1(T^2)\cong\mathbb Z\oplus\mathbb Z$ that bounds a meridian disk in the attached solid torus.  Thus a link surgery has one slope per component, each defined up to sign.
\end{definition}

\begin{theorem}[Surgery presentation of oriented three-manifolds]
\label{thurston-ch01-surgery-presentation}
\group{thurston-ch01-topology}
\source{thurston:page-0010}
\notready
Every oriented three-manifold is obtainable by Dehn surgery on a link in $S^3$.
\end{theorem}
\begin{proof}
\sketch
Represent the oriented three-manifold by surgery on a framed link in $S^3$, equivalently by removing tubular neighborhoods and regluing solid tori with the prescribed primitive slopes.
\end{proof}

\begin{remark}[Homology of surgery slopes]
\label{thurston-ch01-surgery-homology}
\group{thurston-ch01-topology}
\source{thurston:page-0010}
\notready
For surgery on a knot, choose a basis of the boundary-torus homology in which $(1,0)$ maps to a generator of $H_1$ of the complement and $(0,1)$ maps to zero.  Filling along the primitive slope $(p,q)$ then gives first homology $\mathbb Z/p\mathbb Z$, with $p=0$ interpreted as $\mathbb Z$.  In particular, every knot surgery has cyclic first homology, and every slope $(1,n)$ gives an integral homology sphere.  This cyclicity assertion does not extend to surgery on a link with several components.
\end{remark}

\begin{definition}[Branched covering]
\label{thurston-ch01-branched-cover}
\group{thurston-ch01-coverings}
\source{thurston:page-0010}
\notready
Given a finite-sheeted cover of $S^3-L$, compactify it by adjoining circles covering the components of $L$.  The resulting closed manifold $M$ is a branched cover of $S^3$ over $L$, with canonical projection $p:M\to S^3$ that is a local diffeomorphism away from $p^{-1}(L)$.
\end{definition}

\begin{definition}[Cyclic branched cover]
\label{thurston-ch01-cyclic-branched-cover}
\group{thurston-ch01-coverings}
\source{thurston:page-0010}
\notready
For a knot $K$, the $n$-fold cyclic branched cover is obtained from the cover of $S^3-K$ corresponding to
\[
\ker\bigl(\pi_1(S^3-K)\longrightarrow H_1(S^3-K)\cong\mathbb Z\longrightarrow\mathbb Z_n\bigr).
\]
For the trivial knot these covers are $S^3$.
\end{definition}

\begin{theorem}[Three-sheeted branched-cover presentation]
\label{thurston-ch01-three-sheeted-cover}
\group{thurston-ch01-coverings}
\source{thurston:page-0010}
\notready
Every oriented three-manifold is a three-sheeted branched cover of $S^3$ over some knot.  Such a cover need not be regular, so it may have no covering transformations.
\end{theorem}

\begin{theorem}[Heegaard decomposition]
\label{thurston-ch01-heegaard-decomposition}
\group{thurston-ch01-topology}
\source{thurston:page-0011}
\notready
Every three-manifold is obtained by gluing the boundaries of two handlebodies of a common genus.  The collection of gluing maps is generally complicated, and deciding when two maps yield homeomorphic manifolds is difficult.
\end{theorem}

\begin{definition}[Polyhedral face gluing]
\label{thurston-ch01-polyhedral-gluing}
\group{thurston-ch01-polyhedra}
\source{thurston:page-0011}
\notready
Let $P_1,\ldots,P_k$ be compact polyhedral three-balls with an even number of $K$-gonal faces for every $K$.  Pair every two-face with a face of the same combinatorial type by an orientation-reversing cellular identification.  Assume that these pairings give a finite quotient $X$ whose edge links are circles.  Then $X$ is an oriented three-complex and is a manifold away from the images of the vertices.  The link of each quotient vertex is a closed connected oriented surface, and removing open cone neighborhoods of the non-manifold vertices yields a compact three-manifold with boundary.
\end{definition}

\begin{proposition}[Euler criterion for polyhedral quotients]
\label{thurston-ch01-polyhedral-euler-criterion}
\group{thurston-ch01-polyhedra}
\source{thurston:page-0011}
\notready
Under the preceding face-pairing hypotheses, let $L_v$ be the link of each quotient vertex $v$.  The following are equivalent:
\begin{enumerate}
\item $X$ is a closed oriented three-manifold;
\item every $L_v$ has Euler characteristic $2$, equivalently every $L_v$ is a two-sphere;
\item $\chi(X)=0$.
\end{enumerate}
\end{proposition}
\begin{proof}
\sketch
Remove disjoint open cone neighborhoods of all quotient vertices and call the resulting compact oriented three-manifold $M$.  Its boundary is the disjoint union of the surfaces $L_v$.  Doubling $M$ along its boundary gives a closed oriented three-manifold, so
\[
  0=\chi(DM)=2\chi(M)-\chi(\partial M),
  \qquad
  \chi(M)=\frac12\sum_v\chi(L_v).
\]
Recovering $X$ attaches one cone on each $L_v$.  Additivity of Euler characteristic therefore gives
\[
  \chi(X)
  =\chi(M)+\sum_v\bigl(1-\chi(L_v)\bigr)
  =\frac12\sum_v\bigl(2-\chi(L_v)\bigr).
\]
Each $L_v$ is a closed connected oriented surface, hence $\chi(L_v)\leq2$, with equality exactly for the two-sphere.  Thus $\chi(X)=0$ exactly when every vertex link is a sphere, which is exactly the remaining local condition for $X$ to be a manifold.
\end{proof}

\begin{remark}[Polyhedral gluing counts]
\label{thurston-ch01-polyhedral-counts}
\group{thurston-ch01-polyhedra}
\source{thurston:page-0011}
\notready
The number of face-pairing patterns for an octahedron is
\[
\frac{8!}{2^4\,4!}\,3^4=8505,
\]
and the corresponding count for an icosahedron is $38{,}661$ billion before quotienting by symmetries.  Pairings of the sides of a $2n$-gon giving oriented surfaces number $(2n)!/(2^n n!)$; unlike in dimension three, the Euler characteristic completely classifies closed oriented surfaces.
\end{remark}

\begin{example}[Two-tetrahedron figure-eight complement]
\label{thurston-ch01-figure-eight-tetrahedra}
\group{thurston-ch01-examples}
\source{thurston:page-0012,thurston:page-0013,thurston:page-0014}
\notready
Take two tetrahedra and identify their faces in the unique orientation-compatible pattern matching the edge arrows.  The quotient has two tetrahedra, four triangles, two edges, and one vertex, hence Euler characteristic $1$.  The vertex link is a torus; deleting the vertex therefore gives a three-manifold homeomorphic to the complement of the figure-eight knot.
\end{example}

\begin{remark}[Cell decompositions need not be geometric]
\label{thurston-ch01-cell-decomposition-warning}
\group{thurston-ch01-examples}
\source{thurston:page-0014,thurston:page-0015}
\notready
A knot drawn along the one-skeleton of a tetrahedron determines a cell decomposition of its complement, but the complementary three-cells need not be tetrahedra.  Thus a suggestive diagram does not by itself identify the manifold or its cell structure.
\end{remark}
